\section{Method}
We conducted a controlled, within-subjects laboratory experiment to compare the effects of three audio configurations on remote collaborative teams.

\subsection{Experimental Design}
The study used a single-factor, within-subjects design with three levels of audio spatialization:
\begin{enumerate}[label=\lbrack A\arabic*\rbrack]
\item \textbf{Mono Audio:} All auditory signals were mixed into a single mono track and delivered equally to both ears, simulating a standard telephone conference call.
\item \textbf{Stereo Audio:} Audio signals from partners were panned left or right based on their initial layout in the virtual workspace, without accounting for head rotation.
\item \textbf{3D Spatialized Audio:} Audio signals were dynamically filtered in real time using head-related transfer functions (HRTFs) corresponding to the real-time position and orientation of the avatars.
\end{enumerate}

\subsection{Task and Environment}
Triads of participants worked together in a custom virtual workspace built by Vertex Interaction Group. The environment simulated a control room where teams had to solve a complex scheduling problem under time pressure. The task required members to share unique, locally stored information and arrange a series of transport flights on a shared timeline. Effective team communication was essential, as no single participant held all the required parameters.

\subsection{Participants}
We recruited 72 triads ($N = 216$ participants, age range: 18--45, 114 female) from Meridian and Synapse corporate research offices. All participants reported normal hearing and normal or corrected-to-normal vision. Triads were assigned to isolated physical rooms and interacted solely through their virtual avatars.
